%!TEX root = forallxyyc.tex

\chapter{Notas sobre acessibilidade}

Foi dada atenção especial para tornar esta versão HTML tão acessível quanto possível, especialmente para leitores que usam Tecnologia Assistiva (TA), como leitores de tela como o
\href{https://www.nvaccess.org/download/}{NVDA}
e o
\href{https://www.freedomscientific.com/products/software/jaws/}{JAWS}.
Leitores de tela são softwares desenvolvidos para permitir que pessoas com baixa visão ou sem visão acessem e naveguem por páginas da web como este livro. Alguns usuários videntes também preferem, às vezes, usar softwares de conversão de texto em fala (TTS) para ler páginas da web em voz alta. Os softwares TTS têm menos recursos que os leitores de tela desenvolvidos para usuários não videntes e podem não converter todo o conteúdo deste livro em áudio.

Se você estiver usando a versão HTML deste livro com um leitor de tela ou terminal Braille, ou se estiver ajudando um estudante que dependa dessas ferramentas (por exemplo, como professor ou orientador de acessibilidade), e tiver comentários, entre
\href{mailto:rzach@ucalgary.ca}{entre em contato}!

\paragraph{Navegação e estilos} Adotamos o estilo HTML fornecido pelo
\href{https://vlmantova.github.io/bookml/}{BookML},
desenvolvido para o departamento de matemática de Leeds por
\href{https://eps.leeds.ac.uk/maths/staff/4058/dr-vincenzo-l-mantova}{Vincenzo
Mantova}. Ele oferece um menu de navegação recolhível, botões no topo de cada página para selecionar o tamanho da fonte, alternar entre fonte com serifa e sem serifa e escolher entre as opções de exibição preto sobre branco, escuro sobre sépia ou claro sobre escuro.

O corpo da página começa com um link invisível para pular para o início do conteúdo, passando pela barra lateral de navegação e pelos botões de estilo. Os botões aparecem na seguinte ordem:
\begin{itemize}
  \item \textquotedblleft alternar barra lateral\textquotedblright{} para ativar ou desativar a barra de navegação;
  \item \textquotedblleft configurações da fonte\textquotedblright{} para o tamanho da fonte, fonte com serifa ou sem serifa e modos claro, sépia ou escuro;
  \item \textquotedblleft download\textquotedblright{} para os links de download das versões PDF; e
  \item \textquotedblleft informações sobre a barra de ferramentas\textquotedblright{} para uma janela que mostra os atalhos de teclado para a navegação.
\end{itemize}
Os atalhos de teclado são: as setas esquerda e direita para mover-se para a página anterior e a próxima, e `s` para alternar a barra lateral de navegação.

\paragraph{Símbolos e fórmulas}

Os símbolos lógicos e as fórmulas deste livro são convertidos para
\href{https://www.w3.org/Math/}{MathML} e exibidos usando o
\href{https://www.mathjax.org/}{MathJax}. O MathJax oferece recursos adicionais de
\href{https://docs.mathjax.org/en/latest/basic/accessibility.html}{acessibilidade
} para fórmulas, disponíveis no menu de contexto do MathJax.
Com um mouse, você pode clicar com o botão direito em qualquer fórmula para ativá-lo.
Os leitores de tela anunciarão que as fórmulas podem ser ativadas (por exemplo, o NVDA diz \textquotedblleft aplicativo clicável\textquotedblright{} antes de cada fórmula para indicar a presença do menu do MathJax). Se o seu leitor de tela não ler fórmulas como $\forall x(A(x) \land B(x))$, talvez seja necessário ativar a Saída de fala no submenu Acessibilidade do MathJax e recarregar a página.

Eis alguns dos símbolos mais comuns. Diferentes leitores de tela e aplicativos TTS os pronunciam de maneiras diferentes e, às vezes, nem os pronunciam. Depois de cada símbolo, fornecemos a pronúncia usual entre os lógicos.
\begin{itemize}
  \item Letras maiúsculas e minúsculas em itálico, como $A$ (\textquotedblleft A maiúsculo\textquotedblright{}) e $x$ (\textquotedblleft x minúsculo\textquotedblright{})
  \item Letras maiúsculas e minúsculas em escrita caligráfica, como $\metav{A}$ (\textquotedblleft A maiúsculo caligráfico\textquotedblright{}) e $\metav{x}$ (\textquotedblleft x minúsculo caligráfico\textquotedblright{})
  \item Conectivos lógicos: $\enot$ (\textquotedblleft não\textquotedblright{}), $\eor$ (\textquotedblleft ou\textquotedblright{}), $\eand$ (\textquotedblleft e\textquotedblright{}), $\eif$ (\textquotedblleft somente se\textquotedblright{}), $\eiff$ (\textquotedblleft se, e somente se,\textquotedblright{}), $\ered$ (\textquotedblleft contradição\textquotedblright{})
  \item Quantificadores: $\forall$ (\textquotedblleft para todo\textquotedblright{}), $\exists$ (\textquotedblleft existe\textquotedblright{})
  \item Símbolos metateóricos: $\therefore$ (\textquotedblleft portanto\textquotedblright{}), $\proves$ (\textquotedblleft prova\textquotedblright{}), $\nproves$ (\textquotedblleft não prova\textquotedblright{}), $\entails$ (\textquotedblleft acarreta\textquotedblright{}), $\nentails$ (\textquotedblleft não acarreta\textquotedblright{})
  \item Operadores modais: $\ebox$ (\textquotedblleft caixa\textquotedblright{}), $\ediamond$ (\textquotedblleft diamante\textquotedblright{})
\end{itemize}

Os subscritos devem ser pronunciados pelos leitores de tela, embora eles talvez não indiquem que se trata de subscritos. A fórmula $A_2$ pode ser lida como \textquotedblleft A maiúsculo, subscrito 2\textquotedblright{} ou simplesmente como \textquotedblleft A2\textquotedblright{}.

As expressões `\blank{}` e `\ifeff` (`if` com dois `f`) são usadas em todo o livro. Na versão HTML, `\blank{}` inclui a palavra soletrada invisível `blank`, que os leitores de tela devem ler em voz alta. A abreviação `\ifeff` significa «se, e somente se». Na versão HTML, ela é codificada usando um caractere invisível de junção de palavras sem espaço, que deve fazer os leitores de tela pronunciarem `\ifeff` como `if-eff` ou `aye-eff-eff`.

\paragraph{Provas} As provas de dedução natural em
\cref{ch.NDTFL,ch.NDFOL,ch.ML} usam linhas verticais para indicar onde as subprovas começam e terminam. Essas linhas verticais se estendem da linha da suposição da subprova até sua última linha e são exibidas entre os números das linhas e as fórmulas em cada linha. Isso torna as provas especialmente desafiadoras para usuários com baixa visão ou perda completa da visão.

Para tornar essas provas acessíveis nesta versão HTML, elas são codificadas como tabelas. Cada linha da tabela tem quatro colunas: o número da linha, um indicador do nível da subprova, uma fórmula e uma justificativa. O indicador do nível da subprova é um número que registra em quantas subprovas aninhadas a linha atual está contida. Ele é 0 se a linha não estiver contida em uma subprova, 1 se estiver em uma subprova, 2 se estiver em uma subprova aninhada em outra subprova e assim por diante. Ao ler um indicador de nível de subprova, os leitores de tela também devem anunciar se uma subprova acabou de ser fechada na linha anterior e quando começa uma nova subprova. As linhas de cabeçalho das tabelas e os indicadores de nível das subprovas ficam ocultos, de modo que as provas aparecem visualmente como no texto impresso.

Eis um exemplo desse tipo de prova:
\begin{fitchproof}
\hypo{wxyz}{(W \eor X) \eor (Y \eor Z)}\PR
\hypo{xy}{X \eif Y}\PR
\hypo{nz}{\enot Z}\PR
\open
	\hypo{wx}{W \eor X}\AS
	\open
		\hypo{w}{W}\AS
		\have{wy1}{W \eor Y}\oi{w}
	\close
	\open
		\hypo{x}{X}\AS
		\have{y1}{Y}\ce{xy, x}
		\have{wy2}{W \eor Y}\oi{y1}
	\close
	\have{wy3}{W \eor Y}\oe{wx, w-wy1, x-wy2}
\close
\open
	\hypo{yz}{Y \eor Z}\AS
	\have{y}{Y}\ds{yz, nz}
	\have{wy}{W \eor Y}\oi{y}
\close
\have{wy4}{W \eor Y}\oe{wxyz, wx-wy3, yz-wy}
\end{fitchproof}
Ela tem 14 linhas, com 3 premissas, 2 níveis de aninhamento de subprovas e dois pares de subprovas adjacentes. Por exemplo, a subprova que começa na linha 5 termina na linha 6, e a linha 7 inicia outra subprova. Assim, os níveis de subprova das linhas 6 e 7 são iguais, mas as linhas 6 e 7 estão em subprovas diferentes. Se você não consegue ver as linhas das subprovas, precisa prestar atenção especial a como os números dos níveis de subprova mudam e a quando uma fórmula é uma suposição. Um leitor de tela deve ler a linha 7 aproximadamente assim: \textquotedblleft 7, fechar subprova, 2, abrir subprova, X maiúsculo, AS\textquotedblright{}. As abreviações `\PR` (de `Premise`, premissa) e `\AS` (de `Assumption`, suposição) são codificadas na versão HTML usando a tag `\verb|abbr|`.

Observe que aplicativos simples de conversão de texto em fala talvez não leiam em voz alta tabelas complexas, como as que contêm provas ou tabelas-verdade.

\paragraph{Diagramas}

O texto também contém alguns diagramas. Na versão HTML, eles são fornecidos com descrições de imagem, por exemplo:
\begin{center}
\begin{tikzpicture}[alt={Há três formas: a primeira é um círculo cinza identificado como A, há algum espaço vazio, a segunda é um círculo branco identificado como C e a quarta é um quadrado branco identificado como D.}]
	\node[circle, grey_shape] (cat1) {A};
	\node[right=10pt of cat1, diamond, phantom_shape] (cat2)  { } ;
	\node[right=10pt of cat2, circle, white_shape] (cat3)  {C} ;
	\node[right=10pt of cat3, white_shape] (cat4)  {D};
\end{tikzpicture}
\bmlDescription{Há três formas: a primeira é um círculo cinza identificado como A, há algum espaço vazio, a segunda é um círculo branco identificado como C e a quarta é um quadrado branco identificado como D.}
\end{center}
A descrição da imagem deve ser: \textquotedblleft Há três formas: a primeira é um círculo cinza identificada como A, há algum espaço vazio, a segunda é um círculo branco identificada como C e a quarta é um quadrado branco identificada como D.\textquotedblright{}

\paragraph{Tabela de símbolos}

Segue uma tabela com todos os símbolos usados no texto, a maneira mais provável de serem pronunciados pelos leitores de tela e o nome que recebem no texto. Na coluna mais à direita, fornecemos uma forma sugerida de inseri-los usando caracteres ASCII, caso a inserção de símbolos especiais (por exemplo, em um trabalho ou e-mail para o professor) não seja uma opção.

\begin{tabular}{lllll}
  \lxBeginTableHead{}Símbolo\lxTableColumnHead{} & Pronúncia\lxTableColumnHead{} & Significado\lxTableColumnHead{} & Equivalente ASCII\lxTableColumnHead{}\\
  \hline\lxEndTableHead
  $\enot$ & sinal de não & não lógico & \verb+~+ ou \verb+-+\\
  $\eor$ & ou & ou lógico & \verb+\/+\\
  $\eand$ & e & e lógico & \verb+/\+ ou \verb+&+\\
  $\eif$ & seta para a direita & condicional & \verb+->+ ou \verb+>+\\
  $\eiff$ & seta esquerda-direita & bicondicional & \verb+<->+ ou \verb+<>+\\
  $\ered$ & t invertido & contradição & \verb+_|_+ ou \verb+!?+\\
  $\forall$ & para todo & quantificador universal & \verb|A| ou \verb|@|\\
  $\exists$ & existe & quantificador existencial & \verb|E| ou \verb|3|\\
  $\therefore$ & portanto & portanto & \verb|:.|\\
  $\proves$ & t para a direita & prova & \verb+|-+\\
  $\nproves$ & não prova & não prova & \verb+|/-+\\
  $\entails$ & verdadeiro & acarreta & \verb+|=+\\
  $\nentails$ & não verdadeiro & não acarreta & \verb+|/=+\\
  $\ebox$ & quadrado branco & necessário & \verb+[]+\\
  $\ediamond$ & diamante branco & possível & \verb+<>+
\end{tabular}

Os subscritos podem ser representados por um sublinhado; por exemplo, $A_2$ como \verb|A_2|.
